% GENERATED FILE. Edit canonical lessons, then run npm run generate:books.
% canonical-source-hash: fnv1a64:b63a9e5b50e059c5
% canonical-lessons: MW-C166-jagno, MW-C166-khangalno, MW-C166-nid, MW-C166-podhno, MW-C166-sajh

\chapter{Up at Daybreak}
\label{ch:mw-up-at-daybreak}
\hlchaptermodality{166}

\begin{chapteropening}
\textbf{By the end of this chapter:} \emph{I can say when I wake up, rinse my mouth and lie down.}

Complete the last lesson of chapter 166: I can say when I wake up, rinse my mouth and lie down.
\end{chapteropening}

\section[\texorpdfstring{\mw{जागणो} (jāgṇo)}{जागणो (jāgṇo)}]{\mw{जागणो} (jāgṇo) — to wake up}

\label{lesson:MW-C166-jagno}



\begin{quote}
Before the new one: say the Marwadi for brave, then the Marwadi for crazy.
\end{quote}

\subsection*{You'll want to know: \mw{जागणो}}
\textbf{\mw{जागणो}} — \emph{jāgṇo} — \textquotedblleft{}to wake up\textquotedblright{}.

\begin{quote}
\mw{हूं} \mw{रोज} \mw{तड़कै} \mw{जागूं} \mw{हूं।}
\end{quote}

\begin{quote}
\emph{hū̃ roj taṛkai jāgū̃ hū̃.}
\end{quote}

\begin{quote}
I wake up at daybreak every day.
\end{quote}

Past (I …): \mw{हूं} \mw{आज} \mw{तड़कै} \mw{जाग्यो}, \emph{hū̃ āj taṛkai jāgyo}.

\begin{cousinweb}
\textbf{\mw{जागणो}} (\emph{jāgṇo}) is inherited from Sanskrit \emph{jāgrati}, \textquotedblleft{}they are awake\textquotedblright{}; CDIAL lists Marwari \emph{jāgṇo}, \textquotedblleft{}to wake\textquotedblright{}, under it (CDIAL 5175). Its partner \textbf{\mw{जगावणो}}, to wake someone, is already yours.
\end{cousinweb}

\begin{grammarlens}[title={What you've built}]
One more word for saying what you do.
\end{grammarlens}

\subsection*{Your turn}
\begin{itemize}
\raggedright
  \item \emph{Say it:} \emph{jāgṇo}
  \item \emph{Say it:} \emph{jāgṇo}, once more
  \item \emph{Say it:} \emph{hū̃ roj taṛkai jāgū̃ hū̃.} — I wake up at daybreak every day.
  \item \emph{Recall:} say the Marwadi for brave, then the Marwadi for crazy, then say \emph{jāgṇo} again
\end{itemize}

\begin{tcolorbox}[breakable,colback=teal!4,colframe=teal!35!black,title={Before you move on}]
What does \mw{जागणो} mean? (\textquotedblleft{}To wake up\textquotedblright{}.) Say it once more.
\end{tcolorbox}
\section[\texorpdfstring{\mw{खंगाळणो} (khaṅgāḷṇo)}{खंगाळणो (khaṅgāḷṇo)}]{\mw{खंगाळणो} (khaṅgāḷṇo) — to rinse}

\label{lesson:MW-C166-khangalno}



\begin{quote}
Before the new one: say the Marwadi for crazy, then the Marwadi for to wake up.
\end{quote}

\subsection*{You'll want to know: \mw{खंगाळणो}}
\textbf{\mw{खंगाळणो}} — \emph{khaṅgāḷṇo} — \textquotedblleft{}to rinse\textquotedblright{}.

\begin{quote}
\mw{हूं} \mw{दिनूगै} \mw{मूंडो} \mw{खंगाळूं} \mw{हूं।}
\end{quote}

\begin{quote}
\emph{hū̃ dinūgai mū̃ḍo khaṅgāḷū̃ hū̃.}
\end{quote}

\begin{quote}
I rinse my mouth early in the morning.
\end{quote}

Past (I …): \mw{म्हैं} \mw{लोटो} \mw{खंगाळ्यो}, \emph{mhaĩ loṭo khaṅgāḷyo}.

\begin{cousinweb}
\textbf{\mw{खंगाळणो}} (\emph{khaṅgāḷṇo}) comes from a reconstructed \emph{*khakkhālayati}, \textquotedblleft{}rinses\textquotedblright{}, which CDIAL thinks imitates the sound of rinsing; it lists Hindi \emph{khãgālnā} and a Gujarati verb for rinsing the mouth in the morning (CDIAL 3762). CDIAL gives no Marwari form in this entry.
\end{cousinweb}

\begin{grammarlens}[title={What you've built}]
One more word for saying what you do.
\end{grammarlens}

\subsection*{Your turn}
\begin{itemize}
\raggedright
  \item \emph{Say it:} \emph{khaṅgāḷṇo}
  \item \emph{Say it:} \emph{khaṅgāḷṇo}, once more
  \item \emph{Say it:} \emph{hū̃ dinūgai mū̃ḍo khaṅgāḷū̃ hū̃.} — I rinse my mouth early in the morning.
  \item \emph{Recall:} say the Marwadi for crazy, then the Marwadi for to wake up, then say \emph{khaṅgāḷṇo} again
\end{itemize}

\begin{tcolorbox}[breakable,colback=teal!4,colframe=teal!35!black,title={Before you move on}]
What does \mw{खंगाळणो} mean? (\textquotedblleft{}To rinse\textquotedblright{}.) Say it once more.
\end{tcolorbox}
\section[\texorpdfstring{\mw{नींद} (nī̃d)}{नींद (nī̃d)}]{\mw{नींद} (nī̃d) — sleep}

\label{lesson:MW-C166-nid}



\begin{quote}
Before the new one: say the Marwadi for to wake up, then the Marwadi for to rinse.
\end{quote}

\subsection*{You'll want to know: \mw{नींद}}
\textbf{\mw{नींद}} — \emph{nī̃d} — \textquotedblleft{}sleep\textquotedblright{}.

\begin{quote}
\mw{आज} \mw{म्हारी} \mw{नींद} \mw{पूरी} \mw{कोनी।}
\end{quote}

\begin{quote}
\emph{āj mhārī nī̃d pūrī konī.}
\end{quote}

\begin{quote}
I haven't had my full sleep today.
\end{quote}

A feminine noun.

\begin{cousinweb}
\textbf{\mw{नींद}} (\emph{nī̃d}) is inherited from Sanskrit \emph{nidrā}, \textquotedblleft{}sleep\textquotedblright{}; CDIAL lists Old Marwari \emph{nīda} (CDIAL 7200). It is feminine: \textbf{\mw{म्हारी} \mw{नींद}}.
\end{cousinweb}

\begin{grammarlens}[title={What you've built}]
One more word for saying what you do.
\end{grammarlens}

\subsection*{Your turn}
\begin{itemize}
\raggedright
  \item \emph{Say it:} \emph{nī̃d}
  \item \emph{Say it:} \emph{nī̃d}, once more
  \item \emph{Say it:} \emph{āj mhārī nī̃d pūrī konī.} — I haven't had my full sleep today.
  \item \emph{Recall:} say the Marwadi for to wake up, then the Marwadi for to rinse, then say \emph{nī̃d} again
\end{itemize}

\begin{tcolorbox}[breakable,colback=teal!4,colframe=teal!35!black,title={Before you move on}]
What does \mw{नींद} mean? (\textquotedblleft{}Sleep\textquotedblright{}.) Say it once more.
\end{tcolorbox}
\section[\texorpdfstring{\mw{पोढणो} (poḍhṇo)}{पोढणो (poḍhṇo)}]{\mw{पोढणो} (poḍhṇo) — to lie down (to sleep)}

\label{lesson:MW-C166-podhno}



\begin{quote}
Before the new one: say the Marwadi for to rinse, then the Marwadi for a night's sleep.
\end{quote}

\subsection*{You'll want to know: \mw{पोढणो}}
\textbf{\mw{पोढणो}} — \emph{poḍhṇo} — \textquotedblleft{}to lie down (to sleep)\textquotedblright{}.

\begin{quote}
\mw{दादी} \mw{रोज} \mw{आथण} \mw{जल्दी} \mw{पोढै} \mw{है।}
\end{quote}

\begin{quote}
\emph{dādī roj āthaṇ jaldī poḍhai hai.}
\end{quote}

\begin{quote}
Grandmother lies down early every evening.
\end{quote}

Past (he …): \mw{वो} \mw{खाट} \mw{माथै} \mw{पोढ्यो}, \emph{vo khāṭ māthai poḍhyo}.

\begin{cousinweb}
\textbf{\mw{पोढणो}} (\emph{poḍhṇo}) is inherited from Sanskrit \emph{pravardhate}, \textquotedblleft{}grows up\textquotedblright{}; CDIAL explains the later sense \textquotedblleft{}lie down, sleep\textquotedblright{} from an older \textquotedblleft{}stretched out\textquotedblright{}, and lists Old Marwari \emph{poḍhaï}, \textquotedblleft{}lies down\textquotedblright{} (CDIAL 8789).
\end{cousinweb}

\begin{grammarlens}[title={What you've built}]
One more word for saying what you do.
\end{grammarlens}

\subsection*{Your turn}
\begin{itemize}
\raggedright
  \item \emph{Say it:} \emph{poḍhṇo}
  \item \emph{Say it:} \emph{poḍhṇo}, once more
  \item \emph{Say it:} \emph{dādī roj āthaṇ jaldī poḍhai hai.} — Grandmother lies down early every evening.
  \item \emph{Recall:} say the Marwadi for to rinse, then the Marwadi for a night's sleep, then say \emph{poḍhṇo} again
\end{itemize}

\begin{tcolorbox}[breakable,colback=teal!4,colframe=teal!35!black,title={Before you move on}]
What does \mw{पोढणो} mean? (\textquotedblleft{}To lie down (to sleep)\textquotedblright{}.) Say it once more.
\end{tcolorbox}
\section[\texorpdfstring{\mw{सांझ} (sā̃jh)}{सांझ (sā̃jh)}]{\mw{सांझ} (sā̃jh) — dusk}

\label{lesson:MW-C166-sajh}



\begin{quote}
Before the new one: say the Marwadi for a night's sleep, then the Marwadi for to lie down.
\end{quote}

\subsection*{You'll want to know: \mw{सांझ}}
\textbf{\mw{सांझ}} — \emph{sā̃jh} — \textquotedblleft{}dusk\textquotedblright{}.

\begin{quote}
\mw{सांझ} \mw{घणी} \mw{सांत} \mw{है।}
\end{quote}

\begin{quote}
\emph{sā̃jh ghaṇī sānt hai.}
\end{quote}

\begin{quote}
Dusk is very calm.
\end{quote}

A feminine noun.

\begin{cousinweb}
\textbf{\mw{सांझ}} (\emph{sā̃jh}) is inherited from Sanskrit \emph{sandhyā}, \textquotedblleft{}twilight\textquotedblright{}, the joining of day and night; CDIAL lists Hindi \emph{sā̃jh} (CDIAL 12918). CDIAL gives no Marwari form in this entry.
\end{cousinweb}

\begin{grammarlens}[title={What you've built}]
One more word for saying what you do.
\end{grammarlens}

\subsection*{Your turn}
\begin{itemize}
\raggedright
  \item \emph{Say it:} \emph{sā̃jh}
  \item \emph{Say it:} \emph{sā̃jh}, once more
  \item \emph{Say it:} \emph{sā̃jh ghaṇī sānt hai.} — Dusk is very calm.
  \item \emph{Recall:} say the Marwadi for a night's sleep, then the Marwadi for to lie down, then say \emph{sā̃jh} again
\end{itemize}

\begin{tcolorbox}[breakable,colback=teal!4,colframe=teal!35!black,title={Before you move on}]
What does \mw{सांझ} mean? (\textquotedblleft{}Dusk\textquotedblright{}.) Say it once more.
\end{tcolorbox}
